\subsection{Convexity and exact safe scales}
Write $l(z)=1-\tau^2 z$. The normalized prior contribution is
$l(z)(m(z)-s)$, with
\[
 m(z)=\frac{r e^{-Pz}-a h e^{-Qz}}{1+r e^{-Pz}+h e^{-Qz}}.
\]
Multiplying $l(z)(m(z)-s)>u$ by the positive responsibility
denominator gives
\[
 A_u\,r e^{-Pz}> s l+u+[(a+s)l+u]\,h e^{-Qz}.
\]
Division by $A_u>0$ proves~\eqref{eq:update_band}. In
Proposition~\ref{prop:convex_band}, $A_u$ and both ratio numerators
are positive everywhere by their terminal lower bounds.

For $u\ge0$, the ratio
$V_j(z)=(j l+u)/[(1-s)l-u]$ for $j=s$ or $a+s$ satisfies
\[
 (\log V_j)''=\tau^4
 \left[\frac{(1-s)^2}{A_u^2}-\frac{j^2}{(j l+u)^2}\right]\ge0 .
\]
Multiplication by an exponential of an affine function preserves
log convexity. For $u\tau^2>0$ the displayed inequality is strict,
so both terms are strictly convex. Otherwise $U,W$ are positive
constants and the nonzero exponents give strict convexity.

For $u<0$, the derivatives satisfy
\[
 \begin{aligned}
 U'&=-\frac{|u|\tau^2}{A_u^2},&
 U''&=-\frac{2(1-s)|u|\tau^4}{A_u^3},\\
 W'&=(a+1)U',&
 W''&=(a+1)U'' .
 \end{aligned}
\]
Since $A_u$ decreases with $z$, $|U'|\le d_1$,
$|U''|\le d_2$, $U\ge U_L$, and $W\ge W_L$.
The two exponential-weighted second derivatives are bounded below
by their exponential factors times $P^2U_L-2Pd_1-d_2$ and
$(Q-P)^2W_L-(a+1)d_2$, respectively. The second bound drops
the nonnegative term $-2(Q-P)W'$. Condition~\eqref{eq:band_convexity}
makes both strictly positive. Strict convexity and continuity give
the level interval and the minimum test, including endpoint roots.

For the candidate barrier, the posterior mean in the horizontal
coordinate decreases with $t$. Its derivative equals
\[
 \partial_t m=-\frac{d}{v}w_2[a(1-w_2)+w_1]\le0 .
\]
The horizontal step at $s=\xi$ is therefore largest at $t=-\eta_B$.
Own-coordinate monotonicity extends an inward bound on the face
to the entire half-strip, even after a perturbation of size $\delta$.
The lower-face check also extends by vertical monotonicity.
A sufficient lower-face check is
$\gamma[b_2+(1+c_-)\eta_B]\ge\delta$ with $c(v)\ge c_->0$,
because the vertical posterior mean is nonnegative. For example,
the step bound $\gamma(1+c_+)\le1$ with $0<\alpha\le1$ and $c(v)\le c_+$
gives own-coordinate monotonicity via the nonnegative diagonal
of the posterior covariance. Thus, with the lower-face check,
the horizontal inequality at the corner is necessary and sufficient
for invariance for all allowed errors. At a failing scale, choose
the corner as the state and the largest allowed positive horizontal
error to exhibit a violation.

\subsection{Critical mass for endpoint improvement}
\begin{proof}
The inequalities in~\eqref{eq:critical_regime} give a nonpositive
terminal prior contribution and a positive terminal denominator $A_u$
at the high-endpoint target. The high-endpoint level is nonnegative, so
Proposition~\ref{prop:convex_band} applies. Since $J_u(0)=r$,
strict improvement is equivalent to $J_u'(0)<0$. Differentiating
the normalized prior contribution at zero gives
$D(h)/(1+r+h)^2$, with sign opposite to $J_u'(0)$.
The polynomial starts negative and has positive quadratic
coefficient for $\tau>0$. When $\tau=0$ it is increasing affine.
At the positive root $\partial_hD>0$ and
$\partial_Q D=h[a+(1+a)r]>0$. Implicit differentiation gives
$\partial_Q h_{\rm crit}<0$.
\end{proof}
